\documentclass[12pt,a4paper]{ctexart}
\usepackage{geometry}\geometry{margin=2.2cm}
\usepackage{fancyhdr,hyperref,longtable,booktabs,enumitem,xcolor,listings}
\pagestyle{fancy}\fancyhf{}\fancyhead[L]{IOMan 逆向分析报告}\fancyhead[R]{V1.0}\fancyfoot[C]{\thepage}
\title{\textbf{IOMan.exe 逆向分析报告}\\基于 PDB 调试符号的架构还原}
\author{约翰·刘}\date{2026年7月13日}

\begin{document}\maketitle
\section{逆向目标}
\begin{itemize}[leftmargin=*]
  \item 目标: IOMan.exe (299KB, 32-bit MFC C++)
  \item 符号: IOMan.pdb (1.9MB, 完整调试符号)
  \item 方法: PDB 符号提取 + 导入表分析 + 字符串扫描
\end{itemize}

\section{进程架构}

IOMan 采用 MFC Document/View 架构，核心类层次：

\begin{verbatim}
CMainFrame (主窗口)
  ├── CMainManager (总控制器)
  │     ├── CPSpaceDB ×N (pSpace 数据源)
  │     ├── COpcServer ×N (OPC DA 数据源)
  │     ├── CIOChannel ×N (采集通道)
  │     │     └── CIODevice ×N (设备)
  │     │           └── CIOItem ×N (测点/Tag)
  │     └── CIODataBuffer (数据缓冲)
  └── CIOManager (IO线程管理)
\end{verbatim}

\section{核心类分析 (PDB 符号还原)}

\subsection{CMainManager — 总控制器}

\begin{table}[H]\centering
\begin{tabular}{lp{7cm}}
\toprule
\textbf{方法} & \textbf{功能} \\
\midrule
InitByIoProject(VString) & 从 IoProject 初始化 \\
GetpSpaceConfig(handle, IoDataSourcePSPace*) & 获取 pSpace 数据源配置 \\
GetIoServerConfig(handle, IoServerParam*) & 获取 IO 服务器配置 \\
GetChannelConfig(handle, IoChannelParam*) & 获取通道配置 \\
GetDeviceConfig(handle, ...) & 获取设备配置 \\
GetTagLinkConfig(path, index, IoTagLinkParameter*) & 获取 Tag-设备链接 \\
GetTagLinkCount(path) & Tag 链接计数 \\
GetDataSourceCount() & 数据源计数 \\
InitTagID / SaveTagID / LoadTagID & Tag ID 管理 \\
InitOpcTagID & OPC Tag 初始化 \\
NewIoDataServerByProject(IoDataSourcePSPace*) & 创建 pSpace 数据服务 \\
NewIoDataServer(IoDataSourcePSPace*) & 创建数据服务 \\
NewTag(IoTagLinkParameter*) & 创建新 Tag \\
NewChannelByIoProject(IoChannelParameter*) & 创建通道 \\
NewIoProjectByProject(IoServerSystemParameter*) & 创建 IoProject 连接 \\
OnlineInitIoCommitDBTagID & 在线初始化 IOCommit Tag ID \\
\bottomrule
\end{tabular}
\end{table}

\subsection{CMainFrame — 主窗口}

\begin{table}[H]\centering
\begin{tabular}{lp{7cm}}
\toprule
\textbf{方法/成员} & \textbf{功能} \\
\midrule
DealProject(path) & 处理工程文件 \\
ParseIoDataSource(path) & 解析 IO 数据源配置 \\
DealTag(IoTagLinkParameter*) & 处理 Tag 链接 \\
DealIODataSource(IoDataSourcePSPace*) & 处理数据源 \\
SendDataToIoProject(data, length, flags) & 向 IoProject 发送数据 \\
SendGetIoProjectHwnd() & 获取 IoProject 窗口句柄 \\
OnlineUpdateProject & 在线更新工程 \\
m\_hIOProjectHwnd & IoProject 窗口句柄 \\
m\_hIOProjectEvent & IoProject 同步事件 \\
m\_hIOProjectExitEvent & IoProject 退出事件 \\
m\_pProjectCompleteList & 工程完成列表 \\
\bottomrule
\end{tabular}
\end{table}

\subsection{CPSpaceDB — pSpace 数据库接口}

构造函数: \texttt{CPSpaceDB(CMainManager*, IoDataSourcePSPace*)}

\begin{table}[H]\centering
\begin{tabular}{lp{7cm}}
\toprule
\textbf{方法} & \textbf{功能} \\
\midrule
Connect / DiagServerConnect & 连接/诊断连接 pSpace \\
IsConnect / SetConnectSuccess & 连接状态管理 \\
InitTags / InitTagID / InitTagValue & Tag 初始化和解析 \\
IODataBufferToDB & 数据缓冲写入 \\
CreateIODataBufferThread & 创建数据缓冲线程 \\
\bottomrule
\end{tabular}
\end{table}

\subsection{CIODevice — 设备对象}

\begin{table}[H]\centering
\begin{tabular}{lp{7cm}}
\toprule
\textbf{方法} & \textbf{功能} \\
\midrule
InitByIoProject & 从 IoProject 初始化设备 \\
NewItemByProject(IoTagLinkParameter*) & 从工程创建测点 \\
InitLinkItemByIoProject & 初始化 Tag-设备链接 \\
SaveTagID / LoadTagID & 持久化 Tag ID \\
\bottomrule
\end{tabular}
\end{table}

\section{关键数据结构}

\begin{table}[H]\centering
\begin{tabular}{ll}
\toprule
\textbf{结构体} & \textbf{用途} \\
\midrule
IoDataSourcePSPace & pSpace 数据源参数 (服务器/端口/凭据) \\
IoDataSourceOpcServer & OPC DA 数据源参数 (CLSID/ProgID) \\
IoTagLinkParameter & Tag-设备链接参数 \\
IoServerSystemParameter & IO 服务器系统参数 \\
IoChannelParameter & 通道参数 \\
ProjectInfo & 工程信息 \\
\bottomrule
\end{tabular}
\end{table}

\section{Tag ID 缓存文件格式}

\subsection{发现位置}
\texttt{E:\textbackslash IO ServerOnLine\textbackslash run\textbackslash TagID\_IOCommitDB<N>\_<Station>.dat}

\subsection{文件格式}
\begin{verbatim}
[4B: entry_count]
Repeated N times:
  [4B: string_length]
  [string_length bytes: ASCII tag_path]
  [4B: 0xFFFFFFFF (terminator)]
... metadata strings (IP, station name, type)
\end{verbatim}

\subsection{pSpace.dat 实例 (542 bytes, 15 tags)}
\begin{verbatim}
/gscyc/T28V37NODE/JE05C201T28V37ADL
/gscyc/T28V37NODE/JE05C201T28V37ADY
...
Metadata: IP=172.16.51.8, Station=gscyc, Type=IO_StationType
\end{verbatim}

\subsection{pPluse.dat 实例 (123KB)}
\begin{verbatim}
/CY1C7K/G138VS405/JD010707G138VS405ADL
/CY1C7K/G138VS405/JD010707G138VS405ADY
... (数千条)
\end{verbatim}

\section{双命名空间架构}

通过逆向发现 pSpace Tag 存在两套完全独立的命名体系：

\begin{table}[H]\centering
\begin{tabular}{lll}
\toprule
\textbf{属性} & \textbf{层级路径} & \textbf{数字 ID} \\
\midrule
设备类型 & 油田泵油单元 & 保护装置 \\
Tag 命名 & /CY1C7K/G138VS405/... & 5000, 5500, ... \\
访问方式 & psAPI Tag Query & psAPI Real ReadList \\
缓存文件 & TagID\_CY*.dat / pPluse.dat & 动态解析(无缓存) \\
设备 ID 格式 & G138VS405, JD010707... & 02204060100, 02105100097 \\
数量 & 数千 & 14 块确认 \\
\bottomrule
\end{tabular}
\end{table}

\section{数据流完整链路}

\begin{verbatim}
Oracle PROJECT_IODATASOURCE
    ↓ IoProject 读取
IoProject.exe (共享内存 + WM_COPYDATA)
    ↓ InitByIoProject
CMainManager
    ├── GetTagLinkConfig → IoTagLinkParameter
    │       ↓
    │   CIODevice::NewItemByProject → CIOItem (测点)
    │       ↓
    │   CPSpaceDB::InitTags → pSpace Tag IDs
    │       ↓
    │   CIODevice::SaveTagID → TagID_IOCommitDB*.dat (缓存)
    │
    └── NewIoDataServerByProject
            ↓
        CPSpaceDB::Connect → pSpace Server (:8889)
            ↓
        Real ReadList / Subscribe → 实时数据
            ↓
        IODataBufferToDB → IoMonitor → Oracle
\end{verbatim}

\section{关键发现}

\begin{enumerate}[leftmargin=*]
  \item IOMan 不直接从配置文件读取 Tag 映射——它通过 IoProject IPC 获取
  \item Tag ID 缓存文件是运行时生成的，用于加速重启后的初始化
  \item 保护装置使用数字 Tag ID（动态解析），油田设备使用层级路径（缓存到文件）
  \item CMainFrame::ParseIoDataSource 是数据源配置的解析入口
  \item IoTagLinkParameter 结构体包含完整的 Tag→设备→通道映射关系
\end{enumerate}

\end{document}
